\begin{styletable}

Окрестность
    & $\nh{a}$ или $\nh[\d]{a}$
    & \verb|\nh{a}| или \verb|\nh[\d]{a}|
\\

Проколотая окрестность
    & $\nhp{a}$ или $\nhp[\d]{a}$
    & \verb|\nhp{a}| или \verb|\nhp[\d]{a}|
\\

Предел \newline
    & $\lim_{x \to 0}$
    & \verb|\lim_{x \to 0}|
\\

Предел со стрелкой
    & $A \appr{x \to 42} B$
    & \verb|\appr{x \to 42}|
\\

Верхний предел
    & $\uplim_{n \to \infty}$
    & \verb|\uplim_{n \to \infty}|
\\

Нижний предел
    & $\lowlim_{n \to \infty}$\rule[-14pt]{0pt}{0pt}
    & \verb|\lowlim_{n \to \infty}|
\\

Асимпт. эквивалентность
    & $f(x) \sim g(x)$
    & \verb|f(x) \sim g(x)|
\\

<<o-малое>>
    & $\so{x}$
    & \verb|\so{x}|
\\

<<О-большое>>
    & $\bo{x}$
    & \verb|\bo{x}|
\\

\hline
\end{styletable}

\stylenote{
    Команды \texttt{\textbackslash nh} и \texttt{\textbackslash nhp} без явного опционального аргумента обозначают окрестность и проколотую окрестность размера $\e$ для точки, заданной в качестве обязательного аргумента.
    Если необходимо указать другой размер, он передаётся в квадратных скобках первым аргументом.
}